\subsection{IF-12: genuine CM theta kernels need not commute with the fixed root}
The common-root construction at the end of IF-09 is a positive result
for the specified coordinate-exponential test family. It is not a
root-invariance theorem for every theta jet. We give both a completely
specified CM theta counterexample and an existence statement for the
torsion-translation family of an actual nonzero source CM theta vector.
Neither assertion identifies an individual translate with the prescribed
global automorphic class sum.

\paragraph{The source operator and its normalization.}
For a positive Hermitian form $H$, linear in its first argument, put
\[
 U_x f(z)=f(z+x)\exp\{-\pi H(z,x)-\pi H(x,x)/2\}.       \tag{IF-12a}
\]
This is exactly \texttt{Formalism.L2.Th1.H.opU} in
\texttt{Formalism/LocalToGlobal/TorsionIntegrality/}
\texttt{QuadraticSemicharacters/IsSemiChar.lean}, lines 21--22,
and \texttt{TaylorLogMetric.U} in
\texttt{Formalism/LocalToGlobal/KernelLattices/ColumnData.lean},
lines 101--102. The latter uses
$H(z,x)=\operatorname{area}(\Lambda)^{-1}\sum z_i\overline{x_i}$.
We retain the Gaussian in (IF-12a); replacing this operator by bare
translation would change the jets.

For a polynomial $P=\sum_\beta p_\beta Z^\beta$ write
\[
 J_f(P)=\sum_\beta p_\beta\partial^\beta f(0).
                                                               \tag{IF-12b}
\]
These are ordinary derivatives, with no factorial division. This is
the convention of \texttt{ThDistC.derivAt0} and \texttt{jetC}
(\texttt{ThDistC.lean}, lines 56--61). If $P$ is homogeneous of degree $D$ in one
embedding factor, the source normalized functional is
$\chi_F^{-1}c_D^D J_f(P)$. With several embedding factors it is
$\chi_F^{-1}\prod_\tau c_D(\tau)^{D_\tau}J_f(P)$.
For nonzero $\chi_F,c_D(\tau)$ these common factors cannot turn a
nonzero commutator coefficient into zero. This statement concerns the
fixed multidegree, and does not discard factors between different
degrees.

\paragraph{Statement IF-12.}
Use the genuine first-arrow shapes from IF-11:
\[
 \alpha=(18,17,14,11,8,5,2,0),\qquad
 \nu=(28,27,24,21,18,15,12,10).
\]
Let $D_{10}:T_\alpha(\mathbb Z^8)\longrightarrow
T_\nu(\mathbb Z^8)$ be the determinant insertion, and let
\[
 B_f=\sum_\beta J_f(Z^\beta)A_{\nu,\beta}D_{10}:
 T_\alpha(\mathbb C^8)\longrightarrow T_\nu(\mathbb C^{24})
                                                               \tag{IF-12c}
\]
be the exact transported finite kernel of IF-08. Take $F$ to be the
row root $E_{21}$, on both spaces. There is a genuine CM theta vector
$f$ in the source's \texttt{Th1.H.Th} model, given explicitly below,
such that
\[
                 F B_f-B_f F\ne0.                            \tag{IF-12d}
\]
Its period ratio is $i10^{400}$, and its translation point has order
dividing $12$. It is the unitary translate of the source's normalized
product theta \texttt{F0}, not a formal exponential series. It has
the original semicharacter on the common refined period lattice
$12\Lambda^{192}$.

More generally, let $f_0\ne0$ be any entire theta vector for a full
integral polarized lattice $\Lambda\subset\mathbb C^q$, with the
operator (IF-12a), and let $P\ne0$ be any polynomial. For every
integer $\ell\ge2$ there exist $n\ge0$ and
$x\in\ell^{-n}\Lambda$ such that
\[
                         J_{U_xf_0}(P)\ne0.                  \tag{IF-12e}
\]
The translate retains the original semicharacter on
$\ell^n\Lambda$. Applying this to the polynomial in the proof
below gives (IF-12d) in the torsion-translation family of an actual
nonzero source CM theta vector, even if denominators are restricted
to powers of one chosen allowed odd prime. It does not assert failure
of (IF-12d)'s opposite equality for a particular preassigned translate
or global class sum.

\begin{proof}
First isolate an exact coefficient, including the determinant
insertion. Let $t_\alpha$ be the highest tensor and
$t=D_{10}t_\alpha$. In the target, use the first eight coordinates.
The column-height multiplicities of $\nu$ are
\[
          (a_1,\ldots,a_8)=(1,3,3,3,3,3,2,10).
\]
Thus $r=Ft$ is the single tensor obtained by changing $e_1$ to
$e_2$ in the unique $\bigwedge^1$ factor. All factors of height
at least two contain both $e_1,e_2$ and are killed by $F$.
Consequently $r$ is an ambient unit basis vector and $E_{12}r=t$.
The insertion $D_{10}$ commutes with $F$ because a full determinant
has trivial $\mathfrak{sl}_8$ action. With $\Delta_h(Z)$ the leading
$h$-minor of the $24$ by $8$ matrix $Z$, define
\[
 \begin{split}
 Q(Z)&=\prod_{h=2}^{6}\Delta_h(Z)^3\,
                         \Delta_7(Z)^2\Delta_8(Z)^{10},\\
 P(Z)&=(Z_{11}-Z_{22})Q(Z).
 \end{split}                                                  \tag{IF-12f}
\]
The tensor exterior form is the normalized Fock form by IF-04.
Taking the coefficient of $r$ in the commutator on $t_\alpha$
therefore gives the exact identity
\[
       \langle r,(F B_f-B_fF)t_\alpha\rangle=J_f(P).            \tag{IF-12g}
\]
Indeed the first summand has polynomial coefficient $Z_{11}Q$ and
the second has coefficient $Z_{22}Q$. No scalar $c_\nu$ remains:
it has already been included in the normalized Fock/exterior
identification. The functional $\langle r,-\rangle$ is an integral
dual-Weyl test, since $r$ is an integral unit tensor. The polynomial
$P$ has degree $D=155$ and is nonzero.

We construct an explicit genuine theta vector. Put $T=10^{400}$,
$\tau=iT$, $\Lambda_T=\mathbb Z+iT\mathbb Z$, and
\[
 \vartheta_T(z)=\sum_{n\in\mathbb Z}
                   e^{-\pi T n^2}e^{2\pi i n z},\qquad
 g_T(z)=\frac{e^{\pi z^2/(2T)}\vartheta_T(z)}{\vartheta_T(0)}.
                                                               \tag{IF-12h}
\]
Here $\vartheta_T(0)>0$. Let $L_T$ be the period pair with lattice
$\Lambda_T$ and let $\omega_0=-(1+\tau)$. Then
\[
                g_T=\texttt{Th1.F.th0}(L_T,\omega_0/2).
                                                               \tag{IF-12i}
\]
For completeness, the exact source definition is
$\texttt{th0}(L,h)=U_h\operatorname{sgStar}(L)/
(U_h\operatorname{sgStar}(L))(0)$ in
\texttt{EisensteinTranslationProducts/SemiChar.lean}, line 34.
The supplied identity
\texttt{GlobalTheory.SgStarOddTheta.exp\_mul\_oddTheta\_eq\_deriv\_mul\_sgStar}
in \texttt{Development/AbelianVarietyTheory/Proofs/GlobalInterpolation/}
\texttt{GaussianOddJacobiThetaEqDerivMulSgStar.lean}, lines 100--124,
identifies the Gaussian-gauged odd Jacobi theta with
$\operatorname{sgStar}$ times its derivative at zero. Its odd theta
is defined in \texttt{DifferentiableOddTheta.lean}, lines 19--20, by
\[
 e^{\pi i\tau/4+\pi i(z+1/2)}
       \vartheta_T\bigl(z+(\tau+1)/2\bigr).
\]
Substitute $h=-(\tau+1)/2$ into (IF-12a). The linear exponential
in $z$ cancels, leaving a nonzero constant times
$e^{\pi z^2/(2T)}\vartheta_T(z)$. Division by its value at zero
gives (IF-12i), including its scalar. The two hypotheses of that
source identity are supplied by \texttt{Th1.F.thetaLine} in
\texttt{EisensteinTranslationProducts/EllipticOnePole.lean}, line 27,
and \texttt{Th1.F.sgStarBasic} in
\texttt{EisensteinTranslationProducts/QuasiPeriodConstUnique.lean},
lines 21--22. We use these supplied results as cited theorems.
Equivalently, their
quasiperiod law and normalized value at zero give the same identity.

Set $q=24\cdot8=192$ and
$G_T(Z)=\prod_{i,j}g_T(Z_{ij})$. Thus $G_T$ is exactly the source
$\texttt{Th1.H.F0}(L_T,\omega_0,\operatorname{Fin}24\times
\operatorname{Fin}8)$, by
\texttt{QuadraticSemicharacters/ProductThetaMem.lean}, lines 10--11.
Let $a$ be the $24$ by $8$ real matrix with
\[
 a_{11}=1/4,\quad a_{22}=1/3,\quad
 a_{jj}=1/4\ (3\le j\le8),\quad a_{ij}=0\ (i\ne j),
                                                               \tag{IF-12j}
\]
and put $f_T=U_{iTa}G_T$, with
$H(Z,W)=T^{-1}\sum Z_{ij}\overline{W_{ij}}$.
Direct substitution into (IF-12a), without a gauge change, gives
\[
 f_T(Z)=\vartheta_T(0)^{-q}e^{\pi\sum Z_{ij}^2/(2T)}
      \sum_{n\in\mathbb Z^q}
       e^{-\pi T\|n+a\|^2}e^{2\pi i\langle n+a,Z\rangle}.
                                                               \tag{IF-12k}
\]
All these series and their derivatives converge normally on compact
sets: a fixed derivative contributes only a polynomial in $n$,
whereas the Gaussian exponent is negative quadratic in $n$.

This is an allowed analytic theta function in the source's exact
theta model. Its original product theta has semicharacter
$\prod(-1)^{m_{ij}n_{ij}}$ at lattice points $m+iTn$; substitution
in (IF-12h) verifies the law. Since $12a$ is integral, the Heisenberg
commutator for $iTa$ and any vector in $12\Lambda_T^q$ is one.
Consequently $f_T$ has the original semicharacter, here trivial,
on $12\Lambda_T^q$, and belongs to the source's
$\texttt{Th}(H,1,12\Lambda_T^q)$. The latter definition requires
analyticity and the stated translation law, not self-duality of the
refined lattice. Also $iT$ is imaginary quadratic, so the underlying
elliptic curve and its products are CM; $iTa$ is a torsion point of
order dividing $12$. We do not identify this refined lattice with
the restricted self-dual lattice of a particular \texttt{D7.Amb}.

There is a useful exact finite-degree heat formula. Write
$\Delta=\sum_{i,j}\partial_{X_{ij}}^2$. For a homogeneous polynomial
of degree $D$, ordinary differentiation of the Gaussian gives
\[
 \begin{split}
 S_T&:=\vartheta_T(0)^q e^{\pi T\|a\|^2}(2\pi i)^{-D}
                                                    J_{f_T}(P)\\
 &=\sum_{n\in\mathbb Z^q}
    e^{-\pi T(\|n+a\|^2-\|a\|^2)}
       \left(e^{-\Delta/(8\pi T)}P\right)(n+a).
 \end{split}                                                  \tag{IF-12l}
\]
The heat exponential is a finite sum on this polynomial. To verify
its sign and factor, in one variable divide each derivative by
$2\pi i$. Pairing two derivatives on
$e^{\pi z^2/(2T)}$ contributes
$(\pi/T)/(2\pi i)^2=-1/(4\pi T)$; the exponential differential
operator therefore has coefficient $-1/(8\pi T)$. The multinomial
expansion proves (IF-12l) in all variables. This also verifies that
there is no missing multi-index factorial.

We now give an explicit error bound; no numerical evaluation of an
exponentially small theta value is used. Let $M=\|P\|_1$ be its
coefficient $\ell^1$ norm. Expansion of determinants gives
\[
 M\le 2\prod_{h=2}^{6}(h!)^3(7!)^2(8!)^{10}
                  \le2(8!)^{27}<2\cdot10^{135}.                \tag{IF-12m}
\]
At the diagonal matrix (IF-12j),
\[
           P(a)=-4^{-128}3^{-28},\qquad |P(a)|>10^{-156}.
                                                               \tag{IF-12n}
\]
For $R=\max(1,\|x\|_\infty)$ and
$u=qD^2/(8\pi T)$, termwise differentiation gives
\[
 |(e^{-\Delta/(8\pi T)}P)(x)|\le M R^D e^u,
 \quad |(e^{-\Delta/(8\pi T)}P)(a)-P(a)|\le M(e^u-1).
\]
For example, at each differentiation step a coefficient is
multiplied by at most $D^2$ for each of the $q$ summands of $\Delta$;
the degree decreases, so the displayed bound applies to every term.
Here $qD^2=4{,}612{,}800<10^7$, so $u<10^{-393}$ and
\[
                  M(e^u-1)<4\cdot10^{-258}.                   \tag{IF-12o}
\]
For $n\ne0$, using $|a_{ij}|\le1/3$ and integer coordinates gives
\[
 \|n+a\|^2-\|a\|^2\ge\|n\|^2/3,
 \qquad
 (1+\|n\|_\infty)^D\le e^{D\|n\|^2}.
\]
Since $T\ge2D$ and $\pi\ge3$, the absolute tail in (IF-12l) is at
most
\[
 2M\sum_{n\ne0}e^{-T\|n\|^2/2}
 \le16qM e^{-T/2}
 \le128qM/T^2
 <5\cdot10^{-661}.                                            \tag{IF-12p}
\]
To check the middle inequality, set $r=e^{-T/2}$. The
one-coordinate theta sum is at most $1+4r$; thus the sum over
nonzero vectors is at most $(1+4r)^q-1\le8qr$ when $4qr\le1$.
The last bound uses $e^{T/2}\ge T^2/8$. All of these conditions
hold for the specified integer $T$. The bounds (IF-12o)--(IF-12p)
are smaller in total than $|P(a)|/2$, proving $S_T\ne0$ and hence
$J_{f_T}(P)\ne0$. Equations (IF-12g) and (IF-12l) prove the explicit
counterexample. A homogeneous change to CM period coordinates, or
the common source factors in (IF-12b), multiplies this nonzero
coefficient by a nonzero scalar and does not affect the conclusion.

Finally prove (IF-12e), which does not vary the initial CM datum.
For simplicity choose coordinates on $\mathbb C^q$ and write
$H(z,x)=z^{\mathsf t}C\overline x$, with $C$ invertible. If
$d=\deg P$, multiplying $J_{U_x f_0}(P)$ by
$e^{\pi H(x,x)/2}$ gives a polynomial in $\overline x$, of degree
at most $d$, whose coefficients are holomorphic functions of $x$:
\[
 R(x)=\sum_\beta p_\beta\sum_{\gamma\le\beta}
    {\beta\choose\gamma}(-\pi C\overline x)^{\beta-\gamma}
                                      (\partial^\gamma f_0)(x).
                                                               \tag{IF-12q}
\]
Its part of top antiholomorphic degree is
$f_0(x)P_d(-\pi C\overline x)$, where $P_d$ is the leading
homogeneous part of $P$. Since $C$ is invertible and $P_d\ne0$,
choose a multi-index $\eta$ of degree $d$ with nonzero coefficient
$b_\eta$ in $P_d(-\pi C\overline x)$. Wirtinger differentiation
then gives
\[
             \partial_{\overline x}^{\eta}R(x)
                        =\eta!\,b_\eta f_0(x).                \tag{IF-12r}
\]
Every lower antiholomorphic term vanishes, and holomorphic
coefficients have zero antiholomorphic derivatives. Thus $R$ is not
identically zero. Its nonzero set is nonempty and open. A real basis
of the full lattice identifies
$\bigcup_n\ell^{-n}\Lambda$ with a dense subset of
$\mathbb R^{2q}$, so this set contains a point $x$ proving
(IF-12e). If $\lambda\in\ell^n\Lambda$, integrality of
$\operatorname{Im}H$ on $\Lambda$ makes the commutator of
$U_x$ and $U_\lambda$ equal to one. Hence the original theta law
holds on that refined lattice with its original semicharacter.
\end{proof}

\paragraph{Application to the actual eight-factor source family.}
Choose the nonzero source vector
$f_0=B.\texttt{thZero}(8,24)$.
The supplied theorem \texttt{WeilVac.thZero\_ne\_zero}, in
\texttt{Development/AbelianVarietyTheory/Proofs/IwasawaDecomposition/}
\texttt{ExistsWeilOThZeroTriv.lean}, lines 38--42, proves its
nonvanishing under \texttt{B.StdOK}; \texttt{thZero\_mem}
(lines 80--103) places it in the actual CM theta space. Its period
lattice is full: it is built from a nonzero fractional ideal in the
CM field under its CM embeddings, and coordinate tensor powers of
that lattice. The CM embeddings identify its real span with all the
ambient complex coordinates. The source polarization is integral on
this lattice and has positive scalar matrix $\operatorname{area}^{-1}I$.
For a chosen factor $\tau_0$, use
\[
 P_{\mathrm{all}}(Z)=P(Z_{\tau_0})
          \prod_{\tau\ne\tau_0}\prod_{h=1}^{8}
                         \Delta_h(Z_\tau)^{a_h}.               \tag{IF-12s}
\]
It is nonzero and has degree $155$ in each of the eight factors.
The coefficient of the corresponding target tensor in the root
commutator is exactly $J_{U_xf_0}(P_{\mathrm{all}})$.
The last part of the proof therefore supplies a torsion translate
with a nonzero commutator for the actual IF-11 coefficient shapes.
It does not require the CM theta vector to split over embeddings.

For the algebraicity assertion, assume in addition the source
arithmetic hypothesis \texttt{Crv D cc} $\delta$: this consists of
\texttt{IsCurveIdeal} and
\texttt{IsDiscriminantTwelfthRootOfIdealLattice}. The preceding
analytic nonvanishing assertion requires no such arithmetic
hypothesis. Under \texttt{Crv}, the normalized coefficients of
these finitely many jets are algebraic: use
\texttt{thZero\_coefIn} (lines 236--257) and
\texttt{ta\_thMoved} (lines 311--316), in
\texttt{VacuumEntireAndNormalizedJetsAlgebraic.lean}; the latter
combines the former with \texttt{coefIn\_translate} (lines 53--85
of \texttt{CoefInTranslate.lean}). The general
\texttt{TA.transl} (lines 281--290) also transports an already
established \texttt{TA} property under rational translation.
These are properties of the actual normalized CM vector, not a
consequence of the default map \texttt{desc}. A finite number field
contains the finitely many resulting coefficients. Clearing their
denominators gives an integral map with the same nonzero commutator;
no assertion about its Smith factors is needed for this obstruction.

The relation to the actual characteristic kernels is explicit but
one-directional. The source's
\texttt{GlobalTheory.ColNorms.IndPres.xiInd\_expand}, in
\texttt{Proofs/DivisorProofs/CosetIndicatorThetaTorsionTranslates.lean},
lines 102--114, writes each normalized \texttt{xiInd} kernel as a
finite sum of these exact unitary translates. Its
\texttt{pres\_thInd} (lines 182--248) retains the coefficients,
volume factors, character phases, and rational translation points.
Thus their commutators are the same explicit linear combinations of
(IF-12g). This does not say that a nonzero term survives a particular
combination, nor that every torsion translate is an individual
\texttt{thInd} satisfying all fixed global data. In particular,
(IF-12d) refutes automatic common-root invariance of the genuine
theta-translation family; deciding invariance or period valuations
for a prescribed global automorphic block still requires its actual
combination and global invariant lattices from IF-08.
